%%% =========================================================================
%%% The Explanation Capacity Audit — Draft
%%% =========================================================================
\documentclass[11pt]{article}
\usepackage[margin=1in]{geometry}
\usepackage[utf8]{inputenc}
\usepackage[T1]{fontenc}
\usepackage{amsmath,amssymb,booktabs,hyperref,natbib,graphicx,xcolor,microtype}

\title{\Large\textbf{The Explanation Capacity Audit:\\Standard ML Explanation Practice\\Routinely Exceeds the Limits of Explanation}}

\author{%
  Drake Caraker\textsuperscript{1,*},
  Bryan Arnold\textsuperscript{2},
  David Rhoads\textsuperscript{2}
  \\[4pt]
  \textsuperscript{1}\textit{Oura Health, San Francisco, CA}
  \quad
  \textsuperscript{2}\textit{Independent Researchers}
  \\[2pt]
  \textsuperscript{*}\texttt{drakecaraker@gmail.com}
}
\date{}

\begin{document}
\maketitle

\begin{abstract}
The Explanation Capacity Theorem proves that for any system with symmetry
group $G$ acting on a feature space $V$, the maximum number of independent
stable explanatory facts is $C = \dim(V^G)$---the dimension of the
$G$-invariant subspace.  Any claim beyond $C$ has mean squared error equal
to the full norm of the target (the over-explanation penalty).  We audit 149
datasets across 53 scientific domains and find that 75\% exceed their
explanation capacity under standard ML explanation practice: the feature
rankings reported in typical SHAP analyses contain more pairwise claims
than the correlation structure can support.  The directional prediction---within-group
pairwise comparisons are less stable than between-group comparisons---is
confirmed in 86\% of testable datasets (27 significant confirmations,
1 significant reversal; 19 Bonferroni-corrected).  The effect is robust
across model classes (XGBoost, Random Forest, Ridge; 80\% cross-model
agreement) and explanation methods (TreeSHAP $\approx$ native importance).
The single reversal (Dermatology) identifies a principled scope condition:
the prediction holds for substitutable features competing for signal, not
for complementary features co-occurring as disease indicators.  These
results establish that the gap between what ML models can predict and what
can be stably explained about them is not a deficiency of current methods
but a mathematical limit---one that 75\% of standard practice exceeds.
\end{abstract}

\section{Introduction}

A machine learning model trained on gene expression data names TSPAN8 as
the most important biomarker for distinguishing colon from kidney tumours.
Retrain with a different random seed: the model now names CEACAM5, an
immune-evasion marker sharing zero Gene Ontology terms with TSPAN8.  Both
models achieve indistinguishable accuracy~\citep{caraker2026limits}.

This instability is not a software bug.  The Explanation Capacity
Theorem~\citep{caraker2026limits} proves that for any system with
correlated features, there is a hard upper bound---the \emph{explanation
capacity} $C = \dim(V^G)$---on the number of independent feature-level
claims that any method can stably support.  Claims beyond $C$ have mean
squared error equal to the full norm of the target: the best stable
approximation to beyond-capacity content is the zero
vector~\citep{caraker2026limits}.

How often does published practice exceed this bound?  This paper answers
the question empirically.  We audit 149 datasets from 53 scientific
domains---genomics, neuroscience, particle physics, manufacturing,
veterinary medicine, and 48 others---and find that 75\% exceed their
explanation capacity under standard SHAP analysis with stochastic
regularisation.

\section{Methods}

\subsection{The Explanation Capacity}

For a feature space $V$ with $P$ features falling into $g$ correlation
groups under the symmetry group $G$ (identified by SAGE: hierarchical
clustering on $|\text{Spearman } \rho|$ at threshold $\rho^* = 0.70$),
the explanation capacity is $C = g$.  A standard SHAP analysis ranking all
$P$ features makes $P$ positional claims.  The dataset \emph{exceeds
capacity} if $P > C$.

\subsection{Audit Protocol}

For each of 149 datasets (from scikit-learn, OpenML, and PMLB):
\begin{enumerate}
  \item Train $M = 50$ XGBoost models with stochastic regularisation
    (subsample $= 0.8$, colsample\_bytree $= 0.5$, different random seeds)
  \item Compute TreeSHAP importance (mean $|\text{SHAP}|$) for each model
  \item Identify correlation groups via SAGE at $\rho^* = 0.70$
  \item Compute pairwise flip rates: for each feature pair $(i, j)$, the
    fraction of seeds where the minority ranking occurs
  \item Classify each pair as \emph{within-group} (features in the same
    SAGE group; predicted unstable) or \emph{between-group} (different
    groups; predicted stable)
  \item Test: within-group flip rate $>$ between-group flip rate
    (Mann--Whitney $U$, one-sided, $\alpha = 0.005$)
\end{enumerate}

Sensitivity analysis at $\rho^* \in \{0.50, 0.60, 0.70, 0.80, 0.90\}$.
Model-class robustness tested with Random Forest and Ridge on 5
representative datasets.  Method robustness tested with native feature
importance vs.\ TreeSHAP.

\section{Results}

\subsection{Exceedance}

Of 149 datasets, 111 (74.5\%) exceed their explanation capacity at
$\rho^* = 0.70$.  The 38 non-exceeding datasets are predominantly cases
where features are approximately independent ($g \approx P$), confirming
the theorem's scope condition: no Rashomon property $\Rightarrow$ no
exceedance.

Exceedance is robust to threshold choice: 97\% at $\rho^* = 0.50$,
94\% at 0.60, 75\% at 0.70, 70\% at 0.80, 51\% at 0.90.

\subsection{Directional Prediction}

Of 84 datasets with measurable within-group pairs, 72 (85.7\%) show the
predicted direction (within-group flip rate $>$ between-group flip rate).
27 are significant at $p < 0.005$ (19 Bonferroni-corrected at
$p < 4.4 \times 10^{-5}$).  One dataset shows a significant reversal
(Dermatology; see Section~\ref{sec:reversal}).

The 27:1 confirmation-to-reversal ratio at $p < 0.005$ is the headline
result.  Under the null hypothesis of no directional effect, the expected
ratio is 1:1.

\subsection{Effect Sizes}

The strongest confirmations span diverse domains: Satellite Image (remote
sensing, $d = 1.27$), HAR Activity (wearables, $d = 0.91$), Horse Colic
(veterinary, $d = 0.99$), AP Colon Kidney (genomics, $d = 0.90$),
MiniBooNE (particle physics, $d = 0.92$), Vehicle Silhouette (vehicle
engineering, $d = 0.82$), MiceProtein (neuroscience, $d = 0.67$).

Median Cohen's $d$ across all testable datasets: 0.40 (medium effect).

\subsection{$\eta$ Law Degradation}

The within-group flip rate increases monotonically with within-group
correlation strength on synthetic data with known groups ($r = 0.898$,
$p = 0.006$).  The bimodal gap (within $-$ between) also increases with
$\rho$ ($r = 0.864$, $p = 0.012$).  At $\rho = 0.99$ (near-exact
exchangeability), the gap is $+0.31$; at $\rho = 0.50$ (threshold
boundary), groups are undetectable.

This confirms the $\eta$ law's degradation curve: the prediction is
sharpest when features are exactly exchangeable and degrades continuously
with approximate symmetry.

\subsection{Robustness}

\emph{Model class:} XGBoost, Random Forest, and Ridge agree on the
direction of the within/between flip rate difference in 4 of 5 test
datasets (80\%).  The one disagreement (Diabetes) has Ridge showing a tiny
reversal ($d = -0.33$) while XGBoost and RF both confirm.

\emph{Explanation method:} TreeSHAP and native XGBoost feature importance
agree on direction in all 3 tested datasets.  Native importance shows
stronger effects in some cases (Diabetes: $d = 1.81$ native vs.\ $0.94$
TreeSHAP).

\subsection{The Dermatology Reversal}
\label{sec:reversal}

The single significant reversal is Dermatology ($P = 34$, $g = 22$,
within $= 0.039$, between $= 0.137$, $d = -0.80$, reversal
$p = 3.6 \times 10^{-4}$).  The reversal persists across all thresholds
($\rho^* = 0.50$ to $0.90$), confirming it is structural.

The explanation: dermatology features are clinical indicators that
\emph{co-occur} (scaling + erythema indicate the same disease), not
features that \emph{compete} for signal (as correlated gene expression
features do).  The theorem's prediction requires substitutability:
correlated features must be interchangeable with respect to the model.
When correlated features are functionally complementary---used jointly by
the model rather than selected among---within-group stability is
\emph{higher} than between-group stability.

This identifies a principled scope condition for the correlation-as-proxy-for-exchangeability approach: the SAGE grouping correctly identifies
substitutable features (the common case in high-dimensional continuous
data) but misidentifies complementary features (the uncommon case in
clinical indicator data).

\section{Discussion}

The audit establishes three results.

First, \textbf{standard ML explanation practice routinely exceeds the
explanation capacity}.  75\% of 149 datasets---spanning 53 domains from
genomics to particle physics---have more feature-level ranking claims than
the correlation structure can stably support.  This is not a critique of
any particular study; it is a structural property of how SHAP analyses are
reported when features are correlated.

Second, \textbf{the capacity law's directional prediction is empirically
confirmed across diverse domains}.  Within-group comparisons are less
stable than between-group comparisons in 86\% of testable datasets, with
27 significant confirmations and 1 significant reversal (27:1 ratio).  The
effect is robust across model classes and explanation methods.

Third, \textbf{the framework correctly identifies its own scope}.  The
38 non-exceeding datasets are cases where features are independent (no
Rashomon property).  The 1 reversal is a case where correlation indicates
complementarity rather than substitutability.  Both are predicted by the
theory's preconditions.

\paragraph{Implications for practice.}
The audit does not claim that published studies with exceedance are
``wrong.''  Group-level conclusions may be correct even when feature-level
specificity is unsupported.  A study reporting ``TSPAN8 is the most
important gene'' is over-specified; a study reporting ``the invasion
pathway genes are the most important group'' may be perfectly stable.

The recommendation is a \emph{stability audit}: before reporting
feature-level explanations, compute the SAGE groups, the capacity $C$,
and report only the $C$ group-level facts as definitive.  Flag
within-group rankings as provisional.  This costs minutes and prevents
the systematic publication of beyond-capacity claims as stable findings.

\paragraph{Limitations.}
The audit uses correlation as a proxy for exchangeability.  The
Dermatology reversal demonstrates that this proxy fails for complementary
features.  The audit covers only XGBoost with TreeSHAP as the primary
method; robustness checks on RF and Ridge are limited to 5 datasets.
The within-group flip rates (mean 0.22) are below the exact-exchangeability
prediction of 0.50 because the audit datasets have moderate, not exact,
correlation.  The 149 datasets are a convenience sample from sklearn,
OpenML, and PMLB, not a random sample of published ML studies.

\section*{Data and Code Availability}

All code, results, and the pre-registration protocol are available at
\url{https://github.com/DrakeCaraker/universal-explanation-impossibility}.

\bibliographystyle{plainnat}
\bibliography{references}

\end{document}
